%%%%%%%%%%%%%%%%%%%%%%%%%%%
% PREAMBOLO DEL DOCUMENTO %
%%%%%%%%%%%%%%%%%%%%%%%%%%%
\documentclass[a4paper,11pt,oneside,top=3cm,bottom=3cm,left=3.5cm,right=3.5cm,openright,reqno,table]{book}
% openany - fa iniziare i capitoli direttamente nella pagina successiva
% openright - fa iniziare i capitoli nella prima pagina destra disponibile 
% fleqn  - allinea le formule a sinistra anzichè centrarle
% leqno - dispone la numerazione delle formule sulla sinistra o destra
% reqno - dispone la numerazione delle formule sulla destra
%
\usepackage{packages}
% Per non appesantire troppo questo file
% quasi tutti i pacchetti usati sono salvati in packages.sty
%
\linespread{1.5}
% Per avere la parola BOZZA scritta su tutte le pagine
\usepackage[italian,light,first,bottomafter]{draftcopy}
\usepackage{array}
\usepackage{color}
\usepackage{colortbl}
% funziona solo in modalità PS
% Invece per i PDF ho risolto così:
% pdftk tesi.pdf background bozza.pdf output tesi_bozza.pdf
%
%%%%%%%%%%%%%%%%%%%%%%%%%%%%%%%%%
%   DOCUMENTO VERO E PROPRIO    %
%%%%%%%%%%%%%%%%%%%%%%%%%%%%%%%%%
\begin{document}
% FRONTESPIZIO %
\begin{titlepage}
\changepage{}{}{}{-7.5 mm}{}{}{}{}{}
% parametri per cambiare le dimensioni di una singola pagina in ordine:
% {textheight}{textwidth}{evensidemargin}{oddsidemargin}{columnsep}
% {topmargin}{headheight}{headsep}{footskip}
% se voglio centrare la pagina devo mettere bindingoffset/2
% i primi 5 parametri posso usarli con \changetext


\begin{center}
\includegraphics [width=.15\columnwidth, angle=0]{unisa}\\ % height
\vspace{0.5cm}
{\LARGE \scshape Università degli Studi di Salerno}\\
\vspace{0.5cm}
{\Large Dipartimento di Informatica}\\
\vspace{0.1cm}
{\large Corso di Laurea Triennale in Informatica}\\
\vspace{1.5cm}
{\Large \scshape Tesi di Laurea} \\
\vspace{4cm}
{\Huge \bfseries TITOLO TESI} \\
\vspace{5cm}

\begin{minipage}[t]{7cm}
\flushleft
\textsc{Relatore}

Prof. \textbf{Andrea De Lucia} \\
{\small Università degli studi di Salerno} \\[0.25cm]
\end{minipage}
\hfill
\begin{minipage}[t]{7cm}
\flushright
\textsc{Candidato}

\textbf{Nome Cognome} \\
Matricola: 0123456789
\end{minipage}

\vspace{3cm}

%& & \\
%& Candidato & \\
%& \textbf{Fabiano Pecorelli} & \\
{\small Anno Accademico YYYY-YYYY} %\\
%
%
\begin{comment}
\begin{table}[!h]
\centering
\begin{tabular}{c c c} %p{5cm}c
& Tesi di laurea & \\
& \textbf{Fabiano Pecorelli} & \\
& & \\[0.25cm]
Relatore \\
prof. \textbf{Andrea De Lucia} \\
{\small Università degli studi di Salerno} & & {\small Provincia di Salerno}\\
& & \\[0.5cm]
& {\small A.A. 2015-2016} & \\
\end{tabular}
\end{table}
\end{comment}
%
%
\end{center}

\end{titlepage}
%

\frontmatter
% quello che segue è in numerazione romana e i capitoli non verranno numerati
% se non si vuole che compaia il numero di pagina basta usare il comando:
%\nonumber

% RINGRAZIAMENTI %
\begin{titlepage}

\nonumber
\null \vspace {\stretch{1}}
	\begin{flushright}
%	\begin{verse}
\textit{INSERIRE QUI UNA DEDICA O UNA CITAZIONE} \\[5mm]
%	\end{verse}
	\end{flushright}



\end{titlepage}
% SOMMARIO %
\cleardoublepage
%\selectlanguage{italian}
\begin{abstract}

INSERIRE ABSTRACT
\\[1cm]
\end{abstract} 
% INDICI %
\phantomsection
\addcontentsline{toc}{chapter}{Indice}
\tableofcontents
% Il simbolo * serve per evitare che comapaia nell'indice
\clearpage
%\listoffigures
%\clearpage
%\listoftables
% GLOSSARIO
%\cleardoublepage
\phantomsection
\addcontentsline{toc}{chapter}{Glossario}
% per inserire l'elenco dei simboli e degli acronimi nell'indice
\printglossary
% Per stampare il glossario
% per aggiornarlo si deve eseguire da terminale:
% makeindex -s myDoc.ist -t myDoc.alg -o myDoc.acr myDoc.acn
% per inserire una voce nell'elenco:
% \newglossaryentry{voce_etichetta}{name={voce}, description={descrizione}}
% se non compare direttamente nel testo va inizializzata con:
% \glsadd{voce_etichetta}
% oppure se viene richiamata all'interno del testo:
% \gls{voce_etichetta}
% SIMBOLI E NOTAZIONI %
\cleardoublepage
\phantomsection
\addcontentsline{toc}{chapter}{Elenco delle figure}
% per inserire l'elenco dei simboli e degli acronimi nell'indice
%\printglossary[type=\acronymtype,title=Elenco delle figure]
% Per stampare l'elenco dei simboli
\listoffigures
\cleardoublepage
\phantomsection
\addcontentsline{toc}{chapter}{Elenco delle tabelle}
% per inserire l'elenco dei simboli e degli acronimi nell'indice
%\printglossary[type=\acronymtype,title=Elenco delle figure]
% Per stampare l'elenco dei simboli
\listoftables
% per aggiornarlo si deve eseguire da terminale:
% makeindex -s myDoc.ist -t myDoc.glg -o myDoc.gls myDoc.glo
% per inserire una voce nell'elenco:
% \newglossaryentry{voce_etichetta}{name={voce}, description={descrizione}}
% se non compare direttamente nel testo va inizializzata con:
% \glsadd{voce_etichetta}
% oppure se viene richiamata all'interno del testo:
% \gls{voce_etichetta}

\mainmatter
% quello che segue sarà in numerazione araba e i capitoli verranno numerati
%\part{Studio iniziale}
% CAPITOLI
\phantomsection
%\addcontentsline{toc}{chapter}{Introduzione}
\chapter{Introduzione}
\markboth{Introduzione}{}
% [titolo ridotto se non ci dovesse stare] {titolo completo}
 
\section{Motivazioni e Obiettivi} %\label{1sec:scopo}
I recenti progressi nelle tecnologie hardware e di comunicazione, soprattutto lo sviluppo delle architetture su infrastruttura Cloud e del Machine Learning, stanno stravolgendo il modo di interagire tra macchina e uomo. Questo non solo avviene nell'ambito sociale e quindi a scopo di aiutare l'utente finale, ma anche e soprattutto nell'aiuto allo sviluppatore stesso che crea queste tecnologie.

I campi di ricerca del machine learning come il Deep Learning stanno cambiando in modo marcato e velocemente lo scenario con il quale gli sviluppatori si trovano ad interagire. Se prima gli IDE\footnote{Integrated Development Enviromen, ovvero un ambiente di sviluppo software che racchiude in un unica interfaccia grafica tutte le funzionalità necessarie alla creazione di codice per svariati linguaggi di programmazione.} erano semplici editor di testo corredati con qualche aiuto al programmatore, come la semplificazione del debug del codice, la compilazione e scrittura, oggi, anche grazie al machine learning, stanno diventando sempre più un partner indispensabile per lo sviluppo rapido e per diminuire il time-to-market e i bug nel software.

Lo scopo di questa tesi è quello di esplorare queste tecnologie e apprenderne le capacità con il fine di sviluppare un conversational agent per lo smart development che aiuta lo sviluppatore ad ottenere, tramite l'utilizzo di un modello di intelligenza artificiale e dato un input da parte dello sviluppatore, in output il blocco di codice relativo all'input.
Tale software permetterà allo sviluppatore di accedere in maniera quasi del tutto immediata ad una vasta quantità di funzioni.
\section{Risultati}
Lo sviluppo del conversational agent Yumi, ha portato alla creazione di un ottimo tool su cui poter lavorare ulteriormente in studi futuri per poter migliorare la sua implementazione. Si è evidenziato come, nonostante il mancato utilizzo dei modelli di linguaggio di dimensione più grande, come per l'appunto il già rinomato GPT-3, gli output di Yumi sono comunque utilizzabili per lo sviluppo di un'estensione per IDE. Uno svantaggio di avere un tool che non utilizza il training per la creazione del modello, ma la cosine similarity per la ricerca intelligente del codice in output, è quello di avere una velocità minore in caso di dataset di grandi dimensioni fino ad arrivare anche all'impossibilità di utilizzarlo in dataset che superano il milione di righe. Questo è derivato dal fatto che il dataset va comunque caricato tutto in memoria sul server che lo utilizza per poterne confrontare la similiarità del commento inserito con il codice trovato. Il vantaggio è per l'appunto quello di poter utilizzare questo tool con dataset più piccoli e direttamente in locale senza dover appoggiare ad architetture che espongono API dietro le quali lavorano anche multipli calcolatori con alte prestazioni. Il tool può essere riadattato a qualsiasi contesto e non obbligatoriamente utilizzato al contesto del program synthesis in quanto anche in uno scenario in cui vi è bisogno della ricerca di similiarità di testo all'interno di documenti, può essere di grande aiuto e facilmente riadattabile grazie agli strumenti di sviluppo utilizzati. Sia il client che il server sono disponibili pubblicamente sulla repository GitHub al link \texttt{https://github.com/exSnake/yumi}.


\section{Struttura della tesi}
Il seguente lavoro di tesi è incentrato sullo sviluppo di un Conversational Agent per l'aiuto agli sviluppatori nella generazione di codice.

Nel capitolo 2 verrà mostrato lo stato dell'arte dello sviluppo dei modelli di Machine Learning attualmente in uso. In particolare si parlerà dell'anatomia di un Conversational Agent e l'utilizzo degli stessi nello smart development e di come essi stanno cambiando il modo di interagire con gli ambienti di sviluppo.

Verranno descritti i Dataset che sono al momento presenti, i quali vengono utilizzati per la creazione, il training e il testing di questi modelli, ovvero per valutarne la loro efficacia. Verranno illustrati i modelli di machine learning già presenti, utilizzati per la creazione di Bot già esistenti come GitHub Copilot\footnote{Sviluppato da Github, utilizza OpenAI Codex per suggerire intere funzioni in real-time direttamente dall'IDE \url{https://github.com/features/copilot}}, Tabnine\footnote{Sviluppato da Tabnine, si integra perfettamente con molti linguaggi nell'aiuto al completamento del testo https://www.tabnine.com/}, InCoder\footnote{Descritto nel capitolo \ref{sssec:num1}}, ecc.
Verrà fatta una panoramica generale su come sono stati generati e quale è il loro ruolo nell'odierno panorama informatico.
Nel capitolo 3 è illustrato lo sviluppo di YUMI, un conversational agent che permette di ricevere, dato un input, in output il blocco di codice.
Verrà mostrato il processo di creazione del dataset, le scelte architetturali prese e i risultati ottenuti da esse.
\chapter{Stato dell'arte} %\label{1cap:spinta_laterale}
% [titolo ridotto se non ci dovesse stare] {titolo completo}
%

\begin{citazione}
Questo capitolo illustra lo stato dell'arte e i lavori presenti in letteratura sugli aspetti di ricerca trattati nel nostro studio. ECC ECC...
\end{citazione}

\newpage
\chapter{Conclusioni} %\label{1cap:spinta_laterale}
% [titolo ridotto se non ci dovesse stare] {titolo completo}
%

\begin{citazione}
BREVE SPIEGAZIONE CONTENUTO CAPITOLO
\end{citazione}

\newpage

\chapter{Conclusioni} %\label{1cap:spinta_laterale}
% [titolo ridotto se non ci dovesse stare] {titolo completo}
%


\begin{citazione}
	Questo studio si è posto l'obiettivo di creare un conversational agent a partire da quanto già presente in letteratura e semplificandone la gestione del modello.
	
	Si è focalizzato sul linguaggio DART di cui poco è presente in rete al momento della scrittura. Il dataset stesso, che è stato il lavoro più lungo e articolato, può essere riutilizzato in molti ambiti. 
	
	I risultati ottenuti sono piuttosto soddisfacenti e dal tool creato, grazie anche all'adozione di tecniche di clean architecture e del modello di sviluppo TDD si può procedere, sia a migliorare l'interazione con lo sviluppatore, magari integrandolo in un IDE, sia a migliorare le funzionalità del modello, per esempio effettuando il fine-tuning con il dataset generato verso un modello di linguaggio più avanzato come il GPT-3 o i suoi derivati.
	
	In conclusione le strade da poter intraprendere a partire da questo studio sono molte. Possiamo però affermare che il machine learning e l'intelligenza artificiale hanno imboccato una buona direzione attraverso anche l'utilizzo del Deep Learning che sicuramente porterà in futuri piuttosto brevi a risultati di grande rilevanza nella generazione del linguaggio, sia umano (chatbot, linguaggio naturale), sia sintattico (program synthesis, linguaggio di sviluppo). 
	
	Anche se le sfumature del linguaggio sono ancora troppe per riuscire ad ottenere un  alto livello di dettaglio, il rapido avanzamento nei nuovi modelli sta rendendo molto probabile che la prossima grande svolta possa essere dietro l'angolo, con lo sviluppo delle nuove tecnologie sarà possibile raggiungere tale scopo in un futuro non molto lontano.
	
	La quarta rivoluzione industriale è iniziata e l'intelligenza artificiale farà da pioniere quindi è importantissimo incentivare e sponsorizzare la ricerca e gli studi fatti in questo ambito.

\end{citazione}

\newpage

%\part{Impatto ambientale}

\backmatter
\phantomsection
\chapter{Ringraziamenti}
\markboth{Ringraziamenti}{}
% [titolo ridotto se non ci dovesse stare] {titolo completo}
INSERIRE RINGRAZIAMENTI QUI 
\end{document}
